\section{Koszul homology and translation classes}
\label{sec:koszul}

The finite-resolution strands give geometric representatives of
syzygies. The Koszul complex gives a second description that is
particularly effective for constructing classes and counting them.
We first define it for a Goldbach ring of any fixed arity $k\ge2$.
We then return to binary rings, identify the polygon classes in
weight $33$, and construct the common-shift classes used in the
Betti estimates.

\subsection{The Koszul complex and its gradings}
Let $\kk$ be our coefficient field of characteristic $\ell$, where
$\ell=0$ or $\ell$ is prime, and put
$R_{\mathbb Z}=\mathbb Z[x_a\mid a\in\A]$ and
$R=R_{\mathbb Z}\otimes_{\mathbb Z}\kk=\kk[x_a\mid a\in\A]$.
We reserve $\mathbb K$ for the integral Koszul free resolution of
$\mathbb Z$ over $R_{\mathbb Z}$, with augmentation sending every
$x_a$ to zero. For infinitely many variables, $\mathbb K$ is the
filtered direct limit of the finite Koszul complexes, extended to
$R_{\mathbb Z}$; its homology is $\mathbb Z$ in homological degree
zero and zero otherwise.
Its scalar extension
$\mathbb K_{\kk}:=\mathbb K\otimes_{\mathbb Z}\kk$
is the Koszul free resolution of $\kk$ over $R$.
For an $R$-module $M$, put
\[
 K(M)=\mathbb K_{\kk}\otimes_R M,\qquad
 K=K(G)=\mathbb K_{\kk}\otimes_R G.
\]
Thus
\[
 \mathrm{H}_i(K)_{n,d}=\Tor_i^R(G,\kk)_{n,d}.
\]
As for $F_{n,d}$, we write $K_{n,d}$ for the whole complex in
bidegree $(n,d)$ and $K_{i,n,d}$ for its term in homological degree $i$.
The complex $K$ has free $G$-modules as its terms, but generally has
nonzero positive homology, so it is not a resolution of $\kk$ over $G$.
Our Betti numbers are
$\beta_{i,n,d}=\dim_\kk\mathrm{H}_i(K)_{n,d}$.

For a finite subset $\sigma=\{a_1<\cdots<a_i\}\subseteq\A$, write
$\varepsilon_\sigma=\varepsilon_{a_1}\cdots\varepsilon_{a_i}$, with
$\varepsilon_\varnothing=1$. The $R_{\mathbb Z}$-linear differential
on $\mathbb K$ is
\begin{equation}\label{setup:koszul-differential}
 \mathrm{d}(\varepsilon_\sigma)=
 \sum_{j=1}^i(-1)^{j-1}x_{a_j}\varepsilon_{\sigma\setminus\{a_j\}}.
\end{equation}
The same formula gives the differential after passing to coefficient
classes in $G_{\mathbb Z}$ or $G$.

\phantomsection\label{setup:koszul-grading}
Give the generators bidegrees and homological degrees
\[
 \deg x_a=\deg\varepsilon_a=(a,1),\qquad
 |x_a|=0,\qquad |\varepsilon_a|=1,
\]
and extend both degrees additively to products. Thus, for $|\sigma|=i$,
\[
 \deg(x^\alpha\varepsilon_\sigma)
   =\deg(x^\alpha)+\sum_{a\in\sigma}(a,1),
 \qquad |x^\alpha\varepsilon_\sigma|=i.
\]
In particular, the integral strand has the free abelian terms
\[
 \mathbb K_{i,n,d}
 =\bigoplus_{\substack{\sigma\subseteq\A,\ |\sigma|=i\\
          \sum_a\alpha_a=d-i\\
          \sum_a a\alpha_a+\sum_{a\in\sigma}a=n}}
       \mathbb Z\,x^\alpha\varepsilon_\sigma,
\]
where $\alpha$ ranges over nonnegative exponent vectors of finite
support. The coefficient has bidegree
$(n-\sum_{a\in\sigma}a,d-i)$; $d$ is the monomial degree of the whole
term. The differential lowers the homological degree by one and
preserves $(n,d)$.
For example, $x_3^2\varepsilon_5\varepsilon_7$ has bidegree $(18,4)$
and homological degree $2$, while its coefficient has bidegree $(6,2)$.

To compare different fields, take the same binomial relations over
$\mathbb Z$. Let $I_{\mathbb Z}\subseteq R_{\mathbb Z}$ be the ideal
generated by these relations, and put
\[
 G_{\mathbb Z}=R_{\mathbb Z}/I_{\mathbb Z},\qquad
 K_{\mathbb Z}=\mathbb K\otimes_{R_{\mathbb Z}}G_{\mathbb Z}.
\]
Extension of scalars gives
\[
 G=G_{\mathbb Z}\otimes_{\mathbb Z}\kk,\qquad
 K\cong K_{\mathbb Z}\otimes_{\mathbb Z}\kk.
\]
Since $I_{\mathbb Z}$ is bihomogeneous, tensoring the whole complex
with $G_{\mathbb Z}$ preserves the bidegree decomposition:
\[
 K_{\mathbb Z,i,n,d}
 =\bigoplus_{\substack{\sigma\subseteq\A\\|\sigma|=i}}
   (G_{\mathbb Z})_{\,n-\sum_{a\in\sigma}a,\,d-i}
   \,\varepsilon_\sigma.
\]
Components of negative monomial degree are zero. The generators are
the classes $\overline{x^\alpha}\varepsilon_\sigma$ of the same
total bidegree $(n,d)$; the binomial relations identify coefficient
monomials only within a single bidegree.
The component-basis proof of Proposition~\ref{setup:components} works
over $\mathbb Z$, so these terms remain free abelian. Since $\A$ is
bounded below, each strand $K_{\mathbb Z,n,d}$ is a finite complex of finitely generated
free abelian groups. Consequently these complexes can be compared
before choosing a field. Their integral homology and cell models are
used in \cite{CompanionTopology} to determine characteristic-dependent
Betti numbers.

\subsection{Finite strands and stabilization}
The same finiteness gives a direct stabilization statement for
Koszul complexes. Unlike the compatible choices of minimal
resolutions in Section~\ref{setup:resolution-construction}, the maps
here are the canonical maps on coefficient monomials and exterior
generators. Before taking cutoffs, apply the affine translation of
Proposition~\ref{setup:affine-invariance} if necessary so that
$\A\subseteq\mathbb Z_{>0}$, and put $\A_m=\A\cap[1,m]$.
For a fixed arity $k$, define $G^{<m}$ by restricting all
$k$ labels on each side of every defining relation to $\A_m$; define
$G^{\le m}$ by also restricting the common Goldbach degree to at most
$m$. This extends the binary cutoff notation to the chosen arity.

\begin{proposition}[Stabilization in a fixed bidegree]
\label{setup:stabilization}
Let $\A\subseteq\mathbb Z_{>0}$, fix an arity $k\ge2$, put
$b=\min\A$, and let $\A_m=\A\cap[1,m]$.
Denote by $K_{\mathbb Z}^{<m}$ the Koszul
complex on the variables indexed by $\A_m$, with coefficients in
$G_{\mathbb Z}^{<m}$. Sending each
$\overline{x^\alpha}\varepsilon_\sigma$ to the same expression over
$G_{\mathbb Z}$ induces an isomorphism of $\mathbb Z$-complexes
\[
 K_{\mathbb Z,n,d}^{<m}
 \longrightarrow K_{\mathbb Z,n,d}
 \qquad\text{if }d\ge1\text{ and }m\ge n-(d-1)b.
\]
The same assertion holds with coefficient ring
$G_{\mathbb Z}^{\le m}$ whenever $m\ge n$.
\end{proposition}
\begin{proof}
Every monomial $x^\alpha\varepsilon_\sigma$ of bidegree $(n,d)$
contains $d$ labels counting multiplicities. Each label $a$ therefore
satisfies
\[
 a+(d-1)b\le n.
\]
Thus every label is at most $m$, so all generators of this strand
occur in the stated cutoff.
The relations in each exterior summand are generated over $\mathbb Z$
by
\[
 x^\gamma(x^\mu-x^\nu)\varepsilon_\sigma,
\]
where $x^\mu-x^\nu$ is a defining binomial and both terms have total
bidegree $(n,d)$. The same bound puts every label in each such relation
inside the cutoff. The two strands therefore have identical generators
and relations, and the differential has the same formula.
By positivity, the common weight of each relevant defining
binomial is at most $n$, so $m\ge n$ also includes all required relations
in the relation-weight truncation.
\end{proof}

Before this bound is reached, new relations may identify old coefficient
classes, so the canonical maps need not be injective.

\subsection{From polygonal strands to Koszul classes}
We now take $G$ to be the binary Goldbach ring on the positive odd
primes. The following example makes the connection between the two
descriptions of Tor explicit. In particular, its upper quadrilateral
will produce the common-shift class used in the introduction.

\Needspace{6cm}
\begin{example}[The strand $F_{33,3}$ and its Koszul cycles]\label{ex:weight33}
The vertices are the monomials $x^\alpha\in\mathcal M_{33,3}$, one for each
of the eight unordered partitions of $33$ into three odd primes.
The eleven product-star edges form the graph below. Its three
quadrilaterals and one triangle subdivide a heptagon, so their
oriented boundaries form an integral basis of its cycle group.
They give the four minimal generators in homological degree two
and bidegree $(33,3)$ over every field. We label the triangle
$e_{33,3}$ and the left, right, and upper quadrilaterals
$e'_{33,3}$, $e''_{33,3}$, and $e'''_{33,3}$, respectively.

\begin{center}
\begin{tikzpicture}[line join=round,line cap=round,
    dot/.style={circle,fill=black,inner sep=1.5pt},
    vertex/.style={font=\small,inner sep=2pt},
    edge label/.style={font=\small,inner sep=1pt},
    face label/.style={font=\normalsize},
    ed/.style={line width=.7pt,mid arrow={Stealth[length=1.45mm,width=1mm]}}]
  \coordinate (H) at (0,3.65);
  \coordinate (B) at (-3.75,2.05);  \coordinate (E) at (3.75,2.05);
  \coordinate (F) at (-4.55,-.65);  \coordinate (D) at (4.55,-.65);
  \coordinate (G) at (-2.05,-3.35); \coordinate (A) at (2.05,-3.35);
  \coordinate (C) at (0,.10);
  \fill[blue!12] (H)--(B)--(F)--(G)--(A)--(D)--(E)--cycle;
  % Each arrow points toward the quadratic anchor for the fixed factor.
  \draw[ed] (B)--(H); \draw[ed] (E)--(H);
  \draw[ed] (B)--(F); \draw[ed] (D)--(E);
  \draw[ed] (F)--(G); \draw[ed] (A)--(D); \draw[ed] (A)--(G);
  \draw[ed] (B)--(C); \draw[ed] (C)--(E);
  \draw[ed] (C)--(G); \draw[ed] (A)--(C);
  \foreach \v in {A,B,C,D,E,F,G,H} \node[dot] at (\v) {};
  \node[vertex,above] at (H) {$x_{11}^3$};
  \node[vertex,above left] at (B) {$x_3x_{11}x_{19}$};
  \node[vertex,above right] at (E) {$x_5x_{11}x_{17}$};
  \node[vertex,left] at (F) {$x_7^2x_{19}$};
  \node[vertex,right] at (D) {$x_5^2x_{23}$};
  \node[vertex,below left] at (G) {$x_7x_{13}^2$};
  \node[vertex,below right] at (A) {$x_3x_7x_{23}$};
  \node[vertex] at (0,.69) {$x_3x_{13}x_{17}$};
  \draw[line width=.35pt] (0,.49)--(C);
  \node[face label] at (0,-2.58) {$e_{33,3}$};
  \node[face label] at (-2.78,-.72) {$e'_{33,3}$};
  \node[face label] at (2.78,-.72) {$e''_{33,3}$};
  \node[face label] at (0,2.24) {$e'''_{33,3}$};
  \node[edge label,above] at (-1.88,3.02) {$x_{11}e'_{22,2}$};
  \node[edge label,above] at (1.88,3.02) {$x_{11}e_{22,2}$};
  \node[edge label,left] at (-4.37,.69) {$x_{19}e_{14,2}$};
  \node[edge label,right] at (4.37,.69) {$x_5e_{28,2}$};
  \node[edge label,left] at (-3.51,-2.01) {$x_7e_{26,2}$};
  \node[edge label,right] at (3.51,-2.01) {$x_{23}e_{10,2}$};
  \node[edge label,below] at (0,-3.47) {$x_7e'_{26,2}$};
  \path (B)--(C) node[midway,sloped,below=3pt,edge label] {$x_3e_{30,2}$};
  \path (C)--(E) node[midway,sloped,below=3pt,edge label] {$x_{17}e_{16,2}$};
  \path (G)--(C) node[midway,sloped,above=3pt,edge label] {$x_{13}e_{20,2}$};
  \path (C)--(A) node[midway,sloped,above=3pt,edge label] {$x_3e'_{30,2}$};
\end{tikzpicture}
\end{center}
The two edges incident to $x_{11}^3$ form the translated star $x_{11}F_{22,2}$.
At $x_{11}^3$ every binary move fixes $x_{11}$, and the two other
pairs of weight $22$ are $(3,19)$ and $(5,17)$. Thus every polygon
passing through $x_{11}^3$ uses both edges.

The fixed factor on each edge also describes the corresponding class in
$\Tor_2^R(G,\kk)_{33,3}$. Use $K=\mathbb K_{\kk}\otimes_R G$, with
$\mathbb K_{\kk}$ the field Koszul resolution introduced above.
At a vertex $x^\alpha$ of an oriented
polygon, let $p$ be the fixed label of the incoming edge and $q$ that of
the outgoing edge. When $p\ne q$ the corner contributes
\[
 \frac{x^\alpha}{x_p x_q}\,\varepsilon_p\varepsilon_q.
\]
When $p=q$ the corner contributes zero since $\varepsilon_p^2=0$.
To follow the polygon in the Koszul complex, read off the fixed labels
of its successive edges. The corner above takes us from label $p$ to
label $q$: its differential is
$(x^\alpha/x_q)\varepsilon_q-(x^\alpha/x_p)\varepsilon_p$.
Along an edge $x^\alpha\to x^{\alpha'}$ fixing $p$, the two adjacent
corners therefore contribute opposite terms with exterior factor
$\varepsilon_p$. Their coefficients agree in $G$ by the binary relation
$(x^\alpha/x_p)=(x^{\alpha'}/x_p)$. Summing around the polygon gives a cycle.

To identify this cycle with the face class in Tor, write the oriented
boundary of its generator $e$ as $\partial e=\sum_j s_j x_{p_j}e_j$,
where $e_j$ is the quadratic generator on the $j$th edge and
$s_j=\pm1$. Let $w$ be the sum of the corners lifted to the Koszul
complex $\mathbb K_{\kk}$ over $R$. In the total complex
$F\otimes_R\mathbb K_{\kk}$ the chain
\[
 e\otimes1+\sum_j s_j e_j\otimes\varepsilon_{p_j}+1\otimes w
\]
is a cycle. The Koszul differential of the middle term cancels
$\partial e\otimes1$ with the tensor sign, and its remaining differential
cancels $1\otimes\mathrm{d}w$ edge by edge. The two augmentations give $e$ modulo
$\m$ and $w$ modulo $I$, respectively.

With the indicated face labels, the following cycles $z_T,z_L,z_R,z_U$
represent $e_{33,3},e'_{33,3},e''_{33,3},e'''_{33,3}$, respectively,
in Tor. We write each exterior pair in increasing order.
\[
\begin{aligned}
 z_T&=-x_{23}\varepsilon_3\varepsilon_7+x_{17}\varepsilon_3\varepsilon_{13}-x_{13}\varepsilon_7\varepsilon_{13},\\
 z_L&=-x_{11}\varepsilon_3\varepsilon_{19}+x_{17}\varepsilon_3\varepsilon_{13}-x_{13}\varepsilon_7\varepsilon_{13}+x_7\varepsilon_7\varepsilon_{19},\\
 z_R&=-x_7\varepsilon_3\varepsilon_{23}+x_{13}\varepsilon_3\varepsilon_{17}-x_{11}\varepsilon_5\varepsilon_{17}+x_5\varepsilon_5\varepsilon_{23},\\
 z_U&=-x_{19}\varepsilon_3\varepsilon_{11}+x_{13}\varepsilon_3\varepsilon_{17}-x_5\varepsilon_{11}\varepsilon_{17}.
\end{aligned}
\]
Each term $\pm x_r\varepsilon_p\varepsilon_q$ comes from the vertex
$x_rx_px_q$. In $z_T,z_L,z_R$ this gives one term for every vertex
of the corresponding face. At the shared vertex $x_3x_{13}x_{17}$,
for instance, the left face contributes $x_{17}\varepsilon_3\varepsilon_{13}$
and the right face contributes $x_{13}\varepsilon_3\varepsilon_{17}$.
The partition is the same. The outgoing edge fixes $13$ on the left face
and $17$ on the right. On the upper face the corner at $x_{11}^3$ contributes
$x_{11}\varepsilon_{11}^2=0$, since both incident edges fix $11$.
The three terms of $z_U$ come from its other three vertices.

For the purpose of following these cycles, one may suppress the
degree-two vertex $x_{11}^3$ and join $x_3x_{11}x_{19}$ directly to
$x_5x_{11}x_{17}$. The new edge fixes $11$ and uses the binary relation
$x_3x_{19}=x_5x_{17}$. This preserves the cycle group and turns the
upper quadrilateral into a triangle. It is a simplification of the
monomial graph; the full strand $F_{33,3}$ still contains the basis
monomial $x_{11}^3$.

Traverse the triangle from $x_3x_7x_{23}$ through
$x_3x_{13}x_{17}$ and $x_7x_{13}^2$. The fixed labels are $3,13,7$.
At $x_3x_{13}x_{17}$ the fixed label changes from $3$ to $13$,
giving $x_{17}\varepsilon_3\varepsilon_{13}$. The next edge keeps $13$ and replaces
$3+17$ by $7+13$. At $x_7x_{13}^2$ the corner is
$x_{13}\varepsilon_{13}\varepsilon_7$. The terms of the two corner
differentials still carrying $\varepsilon_{13}$ are therefore
\[
 +x_3x_{17}\varepsilon_{13}
 \qquad\text{and}\qquad
 -x_7x_{13}\varepsilon_{13}.
\]
They cancel in $G$ by the very move we traversed. Writing the second
corner in increasing order gives $-x_{13}\varepsilon_7\varepsilon_{13}$.
Both exterior labels are $1$ modulo $3$, so the ordered residue pair
$(1,1)$ looks the same before and after this transposition. Keeping the
prime labels on the two positions records its sign.

There is also a choice in moving a factor between the coefficient and the
exterior pair at a vertex with three distinct labels. For $a<b<c$,
\[
 \mathrm{d}(\varepsilon_a\varepsilon_b\varepsilon_c)
     =x_a\varepsilon_b\varepsilon_c-x_b\varepsilon_a\varepsilon_c+x_c\varepsilon_a\varepsilon_b.
\]
Thus one corner can be replaced by the two other sides of this
Koszul triangle.
Adding $\mathrm{d}(\varepsilon_3\varepsilon_7\varepsilon_{23})-\mathrm{d}(\varepsilon_3\varepsilon_{13}\varepsilon_{17})$ to $z_T$ gives the five-term representative
\[
 -x_{13}\varepsilon_7\varepsilon_{13}+x_3\varepsilon_7\varepsilon_{23}-x_7\varepsilon_3\varepsilon_{23}+x_{13}\varepsilon_3\varepsilon_{17}-x_3\varepsilon_{13}\varepsilon_{17}.
\]
Making the corresponding replacements at $x_3x_{13}x_{17}$ and
$x_3x_{11}x_{19}$ in $z_L$ amounts to subtracting
$\mathrm{d}(\varepsilon_3\varepsilon_{13}\varepsilon_{17})+\mathrm{d}(\varepsilon_3\varepsilon_{11}\varepsilon_{19})$. It gives
\[
 x_7\varepsilon_7\varepsilon_{19}-x_{13}\varepsilon_7\varepsilon_{13}+x_{13}\varepsilon_3\varepsilon_{17}-x_3\varepsilon_{11}\varepsilon_{19}
       -x_3\varepsilon_{13}\varepsilon_{17}-x_{19}\varepsilon_3\varepsilon_{11}.
\]
The last two terms account for the intermediate steps of these swaps.

In this bidegree $\mathrm{d}:K_{3,33,3}\to K_{2,33,3}$ has a separate column $\mathrm{d}(\varepsilon_a\varepsilon_b\varepsilon_c)$ for each
partition with three distinct parts. Consequently no boundary has a term
$x_p\varepsilon_p\varepsilon_q$ with $p\ne q$. The vertices $x_7^2x_{19}$ and $x_5^2x_{23}$
give such coordinates for $z_L$ and $z_R$. The vertex $x_7x_{13}^2$
gives one for both $z_T$ and $z_L$. The class of $z_U$ is nonzero
because on the triple $(3,11,19)$ it uses only the coordinate
$x_{19}\varepsilon_3\varepsilon_{11}$, whereas the only incoming
boundary at this triple has all three coordinates.

These four classes are independent over every field. Read the
coefficients of $x_7\varepsilon_7\varepsilon_{19}$ and $x_5\varepsilon_5\varepsilon_{23}$ to detect $z_L$ and $z_R$,
and then $x_{13}\varepsilon_7\varepsilon_{13}$ to detect $z_T$. The preceding test on the
triple $(3,11,19)$ detects the remaining $z_U$.
\end{example}

For a cubic weight $n$, the Koszul differential gives a different
graph, which we denote by $\Lambda_n$. Its vertices are the
\emph{prime labels} occurring in cubic partitions of $n$, whereas
the vertices of $\Gamma_{n,3}$ and the chosen minimal-move graph are
\emph{cubic monomials}. Join distinct labels $a,b$ in $\Lambda_n$
when $n-a-b$ is an odd prime. This edge represents the basis term
$x_{n-a-b}\varepsilon_a\varepsilon_b\in K_{2,n,3}$, with $a<b$.
If $Q_w$ denotes the common quadratic class of represented weight
$w$, the vertices represent $\eta_a=Q_{n-a}\varepsilon_a\in K_{1,n,3}$,
and
\[
 \mathrm d(x_{n-a-b}\varepsilon_a\varepsilon_b)=\eta_b-\eta_a.
\]
Thus $K_{2,n,3}\to K_{1,n,3}$ is the incidence map of $\Lambda_n$.
Each partition $n=a+b+c$ into three distinct primes also contributes
$\varepsilon_a\varepsilon_b\varepsilon_c\in K_{3,n,3}$, whose boundary
is the triangle on $a,b,c$. Fill these triangles; other triangles
in $\Lambda_n$ need not be boundaries. The graph cycle space modulo
these filled-triangle boundaries is $\Tor_2^R(G,\kk)_{n,3}$.

Let $s_n$ count the distinct-part partitions and $t_n$ the partitions
$n=2p+q$ with $p\ne q$. Write $v_n,c_n$ for the numbers of vertices
and connected components of $\Lambda_n$, including isolated vertices.
There are $3s_n+t_n$ edges, and the $s_n$ incoming triangle boundaries
have disjoint supports and unit entries. Consequently
\[
 \beta_{2,n,3}=(3s_n+t_n-v_n+c_n)-s_n
             =2s_n+t_n-v_n+c_n.
\]
For $n=33$, the prime vertices are $3,5,7,11,13,17,19,23$.
The graph has $15$ edges and cycle rank $8$. Its four filled triangles
come from $(3,7,23)$, $(3,11,19)$, $(3,13,17)$, and $(5,11,17)$,
leaving four independent homology classes. These are precisely the
four displayed cycles obtained from the faces of $F_{33,3}$ by
recording successive fixed labels. In particular, the upper face
gives the prime-label cycle $3\to17\to11\to3$; this triangle is
unfilled because $3+17+11\ne33$.

A partition $2p+q$ with $p\ne q$ detects a class through $x_p\varepsilon_p\varepsilon_q$
exactly when $\{p,q\}$ lies on a graph cycle, since all incoming boundaries
vanish on this coordinate. In weight $19$ the partitions are $(3,3,13)$,
$(3,5,11)$, and $(5,7,7)$. Both repeated-part edges are pendant, and
$\beta_{2,19,3}=0$.



\subsection{Common shifts and independent classes}
\label{sec:translations}
Let $G$ now be the binary Goldbach ring on a bounded-below support
$\A$, and use $K=K(G)$ for its coefficient Koszul complex.
The upper quadrilateral of Example~\ref{ex:weight33} is one instance
of a construction that works in every monomial degree.
Suppose $\sigma=\{a_1<\cdots<a_d\}$ and $\sigma+h$ both lie in the support
$\A$, for a positive integer $h$. The two lists need not be consecutive
terms of an arithmetic progression. For the binary ring consider
\begin{equation}\label{eq:translated-cycle}
 z_{\sigma,h}=\sum_{j=1}^d(-1)^{j-1}x_{a_j+h}
       \varepsilon_{\sigma\setminus\{a_j\}}.
\end{equation}

\Needspace{8\baselineskip}
\begin{theorem}[Independent translation classes]\label{thm:translation}
Fix $h>0$ and $d\ge2$. As $\sigma$ varies over the $d$-element subsets of
$\A\cap(\A-h)$, the chains $z_{\sigma,h}$ give linearly independent classes
over every field. For each $\sigma$ the class has homological degree $d-1$
and bidegree $(n,d)$, where $n=h+\sum_{a\in \sigma}a$. In particular,
\[
 \beta_{d-1,n,d}\ \ge\
 \#\{\sigma\subseteq\A\cap(\A-h): |\sigma|=d,\ h+\textstyle\sum_{a\in \sigma}a=n\}.
\]
These classes remain independent under enlargement of the support.
\end{theorem}
\begin{proof}
For $j<r$, the terms in the differential omitting both exterior
indices have opposite signs and coefficients $x_{a_j+h}x_{a_r}$
and $x_{a_j}x_{a_r+h}$. Their pair sums agree, so their
difference vanishes in $G$.

Suppose that a nonzero finite combination of the displayed cycles
were a boundary. Let $a$ be the largest label in any $\sigma$ with nonzero
coefficient, and choose such a $\sigma$ containing $a$. The coordinate
$x_{a+h}\varepsilon_{\sigma\setminus\{a\}}$ determines $\sigma$ for the fixed
$h$ and therefore has nonzero coefficient. Any preimage has monomial degree
and homological degree $d$, so its coefficients are scalars. The only
column giving the chosen coordinate is $\mathrm{d}(\varepsilon_\tau)$, where $\tau=(\sigma\setminus\{a\})\cup\{a+h\}$.
Since $d\ge2$, this column has another term with $a+h$ in its
exterior part. No selected cycle has such a term, and no other column
can cancel it: its coefficient label and its exterior labels together
determine $\tau$. The variables of $G$ are independent in monomial degree one,
so this is a contradiction. The proof uses only coefficients $\pm1$
and applies in every larger support.
\end{proof}

\subsection{Consequences for projective dimension}
The independence theorem turns repeated additive differences into
lower bounds for the size and length of a resolution.

\begin{corollary}\label{cor:density}
If the positive differences in $\A$ have unbounded multiplicities,
the binary Goldbach ring on $\A$ has infinite projective dimension over
$R$. In fact each positive homological degree of the linear
strand then has infinitely many generators. Both statements apply to
the positive odd primes.
\end{corollary}
\begin{proof}
An overlap of size $r$ at difference $h$ gives $\binom{r}{i+1}$
independent classes in homological degree $i$. For each fixed $i$ this
becomes arbitrarily large as $r$ increases. On the odd primes, put
$\A_m=\{p\le m:p\text{ odd prime}\}$. Its $\binom{|\A_m|}{2}$
unordered pairs have fewer than $m$ distinct positive differences. Some
difference therefore occurs at least $\binom{|\A_m|}{2}/m$ times,
which tends to infinity by the prime number theorem. Alternatively,
the prime progressions of Green and Tao \cite{GreenTao08} give
arbitrarily large overlaps by shifting each progression one step.
\end{proof}

Appendix~\ref{app:scroll-shift} gives a stronger bound for each fixed
shift using a scroll resolution. In Section~\ref{sec:betti-euler-shifts}
we fix the Goldbach degree instead and select independent classes
whose positive shifts may vary.

